\chapter{Analysis of Entropic-BPI for \texorpdfstring{$\beta > 0$}{beta > 0}}
\label{chap:analysis_pos}

This chapter reproduces Appendix~B.1 of \cite{essakine2026tight} (Lemmas~7--11 and the proof
of the sample complexity for $\beta > 0$), decomposed into lemmas that are each a plausible
single Lean declaration, with every constant recomputed. The structure is the one of
\cite{menard2021fast} transported to the exponential space of \cite{fei2021exponential}: a
concentration inequality for the optimal exponential value (Lemma~\ref{lem:concentration_pos})
gives optimism (Lemma~\ref{lem:optimism_pos}); an auxiliary ``ring'' value
$\Zr^t$ bridges the pessimistic value $\Zl^t$ and the value $Z^{\pi^{t+1}}$ of the greedy
policy (Lemmas~\ref{lem:ring_pos}, \ref{lem:certificate_pos}), so that the computable
certificate $\pi^{t+1}\cert^t$ dominates the gap $Z^* - Z^{\pi^{t+1}}$
(Lemmas~\ref{lem:cert_dominates_gap_pos}, \ref{lem:stopping_pos}) and the stopping rule is
sound (Lemma~\ref{lem:pac_at_stopping_pos}); finally the certificate is bounded under the true
kernel (Lemma~\ref{lem:cert_true_recursion_pos}), unrolled along the trajectory of $\pi^{t+1}$
(Lemmas~\ref{lem:cert_unrolled_pos}, \ref{lem:ratio_bound_pos}) and summed over the episodes
(Lemmas~\ref{lem:progress_pos}, \ref{lem:stopping_time_bound_pos}, \ref{lem:pac_pos}).

\section{Standing setting and notation}
\label{sec:pos_setting}

Throughout this chapter: $\beta > 0$; $\mathcal M$ is a finite-horizon MDP
(Definition~\ref{def:episodic_mdp}) with reward function $r$ with values in $[0, 1]$
(Definition~\ref{def:reward_fn}); $\delta \in (0, 1)$ and $\eps > 0$; $(X, Y, \mathrm{out})$ is a run
of Entropic-BPI (Definition~\ref{def:entropic_bpi}) with parameters $(r, \beta, \delta, \eps, s_1)$
in the episode environment of $\mathcal M$ from $s_1$ (Definition~\ref{def:episode_env}), on a
probability space $(\Omega, \mathcal F, \Prob)$, with stopping time $\tau$. For every $t \ge 0$ the
quantities $n_h^t$, $\phat_h^t$, $b_h^t$, $\Ut_h^t$, $\Ul_h^t$, $\Zt_h^t$, $\Zl_h^t$, $\pi^{t+1}$,
$\cert_h^t$, $\Ur_h^t$, $\Zr_h^t$ are those of Definitions~\ref{def:counts}, \ref{def:emp_trans},
\ref{def:bonus}, \ref{def:backups}, \ref{def:greedy}, \ref{def:certificate}, \ref{def:ring}
computed from the history of the first $t$ episodes of the run, and $\rho_h^t$, $\rho^{*t}_h$
are the rate terms of Definition~\ref{def:rate_min}. Inside a statement about a fixed $t$ we
write $\pi := \pi^{t+1}$ and, for a fixed $(h, s, a)$,
\[
  n := n_h^t(s,a), \quad \phat := \phat_h^t(\cdot \mid s, a), \quad b := b_h^t(s,a), \quad
  \alpha := \alpha(n, \delta), \quad \alpha^* := \alpha^*(n, \delta), \quad
  \rho := \rho_h^t(s,a), \quad \rho^* := \rho^{*t}_h(s,a),
\]
and for a function $f : \Ss \to \R$, $\phat f := \sum_{s'} \phat(s' \mid s,a) f(s')$,
$(p_h f)(s,a) := \sum_{s'} p_h(s' \mid s,a) f(s')$, $\Var_{\phat}(f)$ and $\Var_{p_h}(f)(s,a)$
the variances of $f$ under these probability vectors (Lemma~\ref{lem:emp_trans_simplex}).
We also write
\[
  B_h := e^{\beta (H - h)} \quad \text{(the range parameter of step $h$)}, \qquad
  c_h := e^{\beta (H + 1 - h)} \quad \text{(the clip of step $h$)},
\]
so that $B_h = c_{h+1}$, $1 \le B_h \le c_h$ and $B_H = 1$. The ranges used all the time are:
$Z^*_{h+1}, U^*_h \in [1, c_h]$ and $Z^*_{h+1} \in [1, B_h]$ (Lemma~\ref{lem:opt_exp_value_range});
$Z^\pi_{h+1} \in [1, B_h]$ for every policy $\pi$ (Lemma~\ref{lem:exp_value_range});
$1 \le \Zl_{h+1}^t \le \Zt_{h+1}^t \le B_h$ and $1 \le \Ul_h^t \le \Ut_h^t \le c_h$
(Lemma~\ref{lem:backups_range}); $0 \le \pi_{h+1}\cert_{h+1}^t \le c_{h+1} = B_h$
(Lemma~\ref{lem:certificate_range}). The good event is $\evGood^+ = \evKL \cap \evBern(Z^*, B) \cap \evCnt$
with $B = (B_h)_{h \le H}$ (Definition~\ref{def:good_event}); ``on $\evGood^+$'' means ``for every
$\omega \in \evGood^+$'', and the statements are pointwise in $\omega$ (no measure theory is
involved before Lemma~\ref{lem:pac_pos}).

\textbf{Lean remark.} All the objects of this chapter are functions of $\omega$ through the
history $\mathrm{histAt}\ X\ Y\ t\ \omega$ (\texttt{Essakine2026Tight/EV2026/Run.lean}:
\texttt{UtildeAt}, \texttt{ZtildeAt}, \texttt{certAt}, \texttt{ringZAt}, \ldots), and the
events $\evGood^+$ etc.\ are subsets of $\Omega$; a lemma ``on $\evGood^+$, for all $t, h, s, a$,
$\ldots$'' is a Lean statement with the hypothesis \texttt{hω : ω ∈ goodEvent X Y M s₁ β δ} and
the arguments \texttt{t h s a} (\texttt{EV2026/AnalysisPos.lean}, section \texttt{Run}). No hypothesis
that $(X, Y)$ is a run of Entropic-BPI is needed before Lemma~\ref{lem:stopping_time_bound_pos}: all the
quantities are functions of the history, and the lemmas are about these functions. Every lemma is
first proved for a fixed history \texttt{hist} under the pointwise consequences of the good event at
that history (Definition~\ref{def:pos_hist_concentration}), and the run version is a one-line
corollary (Lemma~\ref{lem:pos_hist_concentration}); both names are listed in the \texttt{\textbackslash lean}
tags. The standing hypotheses are \texttt{hM : M.HasRewardFn r}, \texttt{hr : ∀ h s a, r h s a ∈ Set.Icc 0 1},
\texttt{hβ : 0 < β}, \texttt{hδ : 0 < δ} and \texttt{hδ1 : δ ≤ 1} ($\delta = 1$ is allowed), and
$\eps \ge 0$ suffices in Lemma~\ref{lem:pac_at_stopping_pos}. The
step $h$ of the paper is the Lean index $h - 1$; the backward recursions are indexed by the
number $j = H + 1 - h$ of steps to the end, so ``backward induction on $h$ from $H + 1$ down to
$1$'' is an induction on $j$ from $0$ to $H$, and the ranges $B_h = e^{\beta(H-h)}$,
$c_h = e^{\beta(H+1-h)}$ are \texttt{Real.exp (β * (j - 1))} and \texttt{Real.exp (β * j)}.
The pair $(s, a)$ with $n_h^t(s,a) = 0$ (never visited) is a separate case of every
backward induction below (deviation~1 of Chapter~\ref{chap:mdp}): its backups are the clips
$\Ut = c_h$, $\Ul = 1$, its certificate is $c_h$ and its rate terms are $\rho = \rho^* = 1$.

\section{Concentration at a history and sign-independent cores}
\label{sec:pos_cores}

This section records the auxiliary objects of the Lean proofs of this chapter and of
Chapter~\ref{chap:analysis_neg} (\texttt{EV2026/AnalysisCommon.lean}): the pointwise content of the
good event at a history, the concentration lemmas of Chapter~\ref{chap:pre_concentration} for
probability vectors, and the parts of the proofs below that do not depend on the sign of $\beta$.

\begin{definition}[Concentration at a history]\label{def:pos_hist_concentration}
\lean{Essakine2026Tight.HistConcentration, Essakine2026Tight.bernRange}\leanok
\uses{def:event_kl, def:event_bern, def:good_event, def:counts, def:emp_trans, def:rates, def:opt_value}
A history of $t$ episodes satisfies the concentration inequalities (for $\mathcal M$, $\beta$, $\delta$)
if for all $h \le H$ and $(s, a)$, $\KL(\phat_h^t(\cdot \mid s, a) \,\|\, p_h(\cdot \mid s, a)) \le \alpha(n, \delta)/n$
with $n = n_h^t(s, a)$ (the threshold being $\infty$ for $n = 0$), and for every visited pair
$\abs{(\phat_h^t - p_h)Z^*_{h+1}(s, a)} \le \sqrt{2\Var_{p_h}(Z^*_{h+1})(s, a)\alpha^*(n, \delta)/n} + 3B_h\alpha^*(n, \delta)/n$,
where $B_h = e^{\beta(H - h)}$ if $\beta > 0$ and $B_h = 1 - e^{\beta(H - h)}$ otherwise.
\end{definition}

\begin{lemma}[The good event gives concentration at every history]\label{lem:pos_hist_concentration}
\lean{Essakine2026Tight.histConcentration_histAt, Essakine2026Tight.HistConcentration.klDiv_le_ofReal}\leanok
\uses{def:pos_hist_concentration, def:good_event}
For $\omega \in \evGood^\pm$ and every $t$, the history of the first $t$ episodes satisfies the
concentration inequalities; at a visited pair the KL threshold is $\alpha(n, \delta)/n$.
\end{lemma}

\begin{proof}
\uses{def:good_event, def:event_kl, def:event_bern}
\leanok
These are the first two components of the good event, read at the episode $t$.
\end{proof}

\begin{lemma}[KL inequalities for probability vectors]\label{lem:pos_vec_kl}
\lean{Essakine2026Tight.integral_eq_vecExp_of_measureReal, Essakine2026Tight.variance_eq_vecVar_of_measureReal, Essakine2026Tight.abs_vecExp_sub_le_of_klDiv_le, Essakine2026Tight.vecVar_le_two_mul_vecVar_add_of_klDiv_le, Essakine2026Tight.vecVar_le_two_mul_vecVar_add_of_klDiv_le', Essakine2026Tight.vecVar_le_two_mul_vecVar_add, Essakine2026Tight.vecVar_le_mul_vecExp}\leanok
\uses{lem:kl_bernstein, lem:kl_transport_var, lem:transport_var, lem:var_le_range_mean, lem:pconc_weighted_measure}
Let $p$ be a probability vector on a finite set, $\nu$ a probability measure on it with vector
$q = (\nu\{x\})_x$, $a \ge 0$, $b \ge 0$, $c \in \R$, and $f, g$ with values in $[c, c + b]$. If
$\KL(\sum_x p(x)\delta_x \,\|\, \nu) \le a$, then
$\abs{pf - qf} \le \sqrt{2\Var_q(f)a} + \frac23 ba$, $\Var_q(f) \le 2\Var_p(f) + 4b^2a$ and
$\Var_p(f) \le 2\Var_q(f) + 4b^2a$. Without the KL assumption,
$\Var_p(f) \le 2\Var_p(g) + 2b\,p\abs{f - g}$, and $\Var_p(f) \le b\,pf$ if $c = 0$. (The integral and
the variance under $\nu$ are $qf$ and $\Var_q(f)$.)
\end{lemma}

\begin{proof}
\uses{lem:kl_bernstein, lem:kl_transport_var, lem:transport_var, lem:var_le_range_mean, lem:pconc_weighted_measure}
\leanok
If $b = 0$, $f$ is constant and all the claims are trivial. Otherwise apply
Lemmas~\ref{lem:kl_bernstein}, \ref{lem:kl_transport_var}, \ref{lem:transport_var}
and~\ref{lem:var_le_range_mean} to $f - c$ and $g - c$ (values in $[0, b]$), under the weighted
measure of $p$ (Lemma~\ref{lem:pconc_weighted_measure}) and $\nu$; the means shift by $c$ and the
variances are invariant. On a finite set, $\int f\,d\nu = \sum_x\nu\{x\}f(x)$
(\texttt{integral\_fintype}) and $\Var_\nu(f) = \sum_x \nu\{x\}(f(x) - \nu f)^2$.
\end{proof}

\begin{lemma}[Core of Lemmas~7 and~12]\label{lem:pos_core_concentration}
\lean{Essakine2026Tight.abs_vecExp_sub_le_bonus_add}\leanok
\uses{lem:pos_vec_kl}
Let $p, \nu, q$ be as in Lemma~\ref{lem:pos_vec_kl}, $0 \le \alpha^\star \le \alpha$ (the rates divided by
$n$), $B \ge 0$, $H \ge 1$, and $f, g$ with values in $[c, c + B]$. If $\KL(p \,\|\, \nu) \le \alpha$ and
$\abs{pf - qf} \le \sqrt{2\Var_q(f)\alpha^\star} + 3B\alpha^\star$, then
\[
  \abs{pf - qf} \le 2\sqrt2\sqrt{\Var_p(g)\alpha^\star} + 5B\alpha + 4HB\alpha + \frac1H p\abs{g - f} .
\]
\end{lemma}

\begin{proof}
\uses{lem:pos_vec_kl, lem:sqrt_mul_le}
\leanok
The computation of the proof of Lemma~\ref{lem:concentration_pos}: with $D := p\abs{g - f}$,
$2\Var_q(f)\alpha^\star \le 8\Var_p(g)\alpha^\star + 8BD\alpha^\star + 8B^2\alpha\alpha^\star$, and this is at most
$a_1^2 + a_2^2 + a_3^2 \le (a_1 + a_2 + a_3)^2$ for $a_1 := 2\sqrt2\sqrt{\Var_p(g)\alpha^\star}$,
$a_2 := D/H + 2HB\alpha^\star$ (as $a_2^2 \ge 4(D/H)(2HB\alpha^\star)$) and $a_3 := 2\sqrt2 B\alpha$ (as
$\alpha^\star \le \alpha$). Hence $\abs{pf - qf} \le a_1 + D/H + (2H + 2\sqrt2 + 3)B\alpha$, and
$2\sqrt2 + 3 \le 5 + 2H$.
\end{proof}

\begin{lemma}[Core of the third term of Lemmas~11 and~16]\label{lem:pos_core_third_term}
\lean{Essakine2026Tight.vecExp_sub_vecExp_le_of_klDiv_le}\leanok
\uses{lem:pos_vec_kl}
Let $p, \nu, q$ be as in Lemma~\ref{lem:pos_vec_kl}, $\alpha \ge 0$, $B \ge 0$, $H \ge 1$, and $f$ with
values in $[0, B]$. If $\KL(p \,\|\, \nu) \le \alpha$, then $qf - pf \le \frac1H pf + (5 + 4H)B\alpha$.
\end{lemma}

\begin{proof}
\uses{lem:pos_vec_kl, lem:sqrt_mul_le}
\leanok
As in the proof of Lemma~\ref{lem:certificate_pos}: $2\Var_q(f)\alpha \le 4B\alpha\,pf + 8B^2\alpha^2$ by
Lemma~\ref{lem:pos_vec_kl}, which is at most $(pf/H + HB\alpha)^2 + (2\sqrt2B\alpha)^2$, so
$qf - pf \le pf/H + (H + 2\sqrt2 + \frac23)B\alpha \le pf/H + (5 + 4H)B\alpha$.
\end{proof}

\begin{lemma}[Cores of the certificate recursion]\label{lem:pos_core_recursion}
\lean{Essakine2026Tight.vecExp_le_one_add_mul_vecExp_add, Essakine2026Tight.bonusTerm_le_of_klDiv_le, Essakine2026Tight.three_mul_add_le}\leanok
\uses{lem:pos_vec_kl}
Let $p, \nu, q$ be as in Lemma~\ref{lem:pos_vec_kl}, $0 \le \alpha^\star \le \alpha$, $B \ge 0$, $H \ge 1$ and
$\KL(p \,\|\, \nu) \le \alpha$.
\begin{enumerate}
  \item[(i)] If $g$ has values in $[0, B]$, then $pg \le (1 + \frac1H)qg + 2HB\alpha$.
  \item[(ii)] If $G, Z$ have values in $[c, c + B]$ and $q\abs{G - Z} \le \Delta$, then
        $2\sqrt2\sqrt{\Var_p(G)\alpha^\star} + 5B\alpha + 4HB\alpha \le 6\sqrt{\Var_q(Z)\alpha^\star} + \frac1H\Delta + 23HB\alpha$.
  \item[(iii)] (Arithmetic.) If $\sigma, \Delta, \rho, B' \ge 0$, $B, B' \le C$, $b \le 6\sigma + \frac1H\Delta + 23HB\rho$
        and $\hat g \le (1 + \frac1H)\Delta + 2HB'\rho$, then
        $3b + (1 + \frac3H)\hat g \le 36\sigma + (1 + \frac{13}H)\Delta + 84HC\rho$.
\end{enumerate}
\end{lemma}

\begin{proof}
\uses{lem:pos_vec_kl, lem:sqrt_mul_le}
\leanok
(i) and (ii) are steps (i) and (ii) of the proof of Lemma~\ref{lem:cert_true_recursion_pos}, with every
square root bounded by comparing squares as in Lemma~\ref{lem:pos_core_concentration}. (iii): $(1 + \frac3H)(1 + \frac1H) \le 1 + \frac7H$,
$(1 + \frac3H)2HB'\rho \le 8HB'\rho$ and $69HB\rho + 8HB'\rho \le 77HC\rho \le 84HC\rho$.
\end{proof}

\section{Concentration and optimism}
\label{sec:pos_optimism}

\begin{lemma}[Concentration of the optimal exponential value; Lemma~7 of the paper]\label{lem:concentration_pos}
\lean{Essakine2026Tight.abs_vecExp_sub_le_bonus_add_of_pos, Essakine2026Tight.abs_vecExp_sub_le_bonusAt_add_of_pos}\leanok
\uses{def:good_event, def:backups, def:emp_trans, def:counts, def:rates, def:opt_value, def:entropic_bpi}
On $\evGood^+$, for all $t \ge 0$, $h \in \{1, \dots, H\}$ and $(s,a)$ with $n = n_h^t(s,a) \ge 1$,
\[
  \abs{(p_h - \phat) Z^*_{h+1}(s,a)}
  \le 2\sqrt2 \sqrt{\Var_{\phat}(\Zt^t_{h+1}) \frac{\alpha^*}{n}}
  + 5 B_h \frac{\alpha}{n} + 4 H B_h \frac{\alpha}{n}
  + \frac1H \phat \abs{\Zt^t_{h+1} - Z^*_{h+1}} .
\]
The sum of the first three terms of the right-hand side is the bonus $b = b_h^t(s,a)$
(Definition~\ref{def:bonus}, with the third term $4 H B_h \alpha(n,\delta)/n$; see
Remark~\ref{rem:pos_constants_bonus}), so the conclusion reads
$\abs{(p_h - \phat) Z^*_{h+1}(s,a)} \le b + \frac1H \phat\abs{\Zt^t_{h+1} - Z^*_{h+1}}$.
\end{lemma}

\begin{proof}
\uses{def:good_event, def:event_bern, def:event_kl, lem:opt_exp_value_range, lem:backups_range, lem:kl_transport_var, lem:transport_var, lem:sqrt_add_le, lem:sqrt_mul_le, lem:rates_props, lem:emp_trans_simplex, lem:pos_core_concentration, lem:pos_hist_concentration}
\leanok
Fix $\omega \in \evGood^+$, $t$, $h$, $(s,a)$ with $n \ge 1$. Write $f := Z^*_{h+1} - 1$ and
$g := \Zt^t_{h+1} - 1$; both take values in $[0, B_h - 1] \subseteq [0, B_h]$
(Lemmas~\ref{lem:opt_exp_value_range} and~\ref{lem:backups_range}) and variances are invariant
under the shift by $1$.

\emph{Bernstein event.} Since $\omega \in \evBern(Z^*, B)$ with $B_h \ge \max Z^*_{h+1} - \min Z^*_{h+1}$
and $n \ge 1$ (Definitions~\ref{def:event_bern}, \ref{def:good_event}),
\[
  \abs{(p_h - \phat) Z^*_{h+1}(s,a)}
  \le \sqrt{2 \Var_{p_h}(Z^*_{h+1})(s,a) \frac{\alpha^*}{n}} + 3 B_h \frac{\alpha^*}{n} .
\]

\emph{Transport of the variance.} Since $\omega \in \evKL$, $\KL(\phat \,\|\, p_h(\cdot \mid s,a)) \le \alpha/n$
(Definition~\ref{def:event_kl}), so Lemma~\ref{lem:kl_transport_var} (first inequality, with
$p = \phat$, $q = p_h(\cdot \mid s,a)$, $f$, $b = B_h$) gives
$\Var_{p_h}(Z^*_{h+1}) \le 2 \Var_{\phat}(Z^*_{h+1}) + 4 B_h^2 \alpha/n$, and
Lemma~\ref{lem:transport_var} (first part, with $f$, $g$, $b = B_h$) gives
$\Var_{\phat}(Z^*_{h+1}) \le 2\Var_{\phat}(\Zt^t_{h+1}) + 2 B_h \phat\abs{Z^*_{h+1} - \Zt^t_{h+1}}$.
Hence
\[
  2 \Var_{p_h}(Z^*_{h+1}) \frac{\alpha^*}{n}
  \le 8 \Var_{\phat}(\Zt^t_{h+1}) \frac{\alpha^*}{n}
  + 8 B_h \, \phat\abs{\Zt^t_{h+1} - Z^*_{h+1}} \, \frac{\alpha^*}{n}
  + 8 B_h^2 \frac{\alpha}{n}\frac{\alpha^*}{n} .
\]

\emph{Square roots.} By Lemma~\ref{lem:sqrt_add_le} (twice), the square root of the right-hand
side is at most the sum of the three square roots. The first is
$2\sqrt2\sqrt{\Var_{\phat}(\Zt^t_{h+1})\alpha^*/n}$. For the second, Lemma~\ref{lem:sqrt_mul_le}
with $x = \phat\abs{\Zt^t_{h+1} - Z^*_{h+1}}$, $y = 8 B_h \alpha^*/n$ and $\eta = 2/H$ gives
\[
  \sqrt{x y} \le \frac{\eta x}{2} + \frac{y}{2\eta}
  = \frac1H \phat\abs{\Zt^t_{h+1} - Z^*_{h+1}} + 2 H B_h \frac{\alpha^*}{n} .
\]
For the third, Lemma~\ref{lem:sqrt_mul_le} with $\eta = 1$ gives
$\sqrt{\alpha \alpha^*} \le (\alpha + \alpha^*)/2$, so
$\sqrt{8 B_h^2 \alpha\alpha^*/n^2} \le \sqrt2 B_h (\alpha + \alpha^*)/n$.

\emph{Collecting.} Altogether, using $\alpha^* \le \alpha$ (Lemma~\ref{lem:rates_props}) and
$H \ge 1$,
\begin{align*}
  \abs{(p_h - \phat) Z^*_{h+1}(s,a)}
  &\le 2\sqrt2\sqrt{\Var_{\phat}(\Zt^t_{h+1})\frac{\alpha^*}{n}}
     + \frac1H \phat\abs{\Zt^t_{h+1} - Z^*_{h+1}}
     + \big(\sqrt2 \alpha + 3 \alpha^*\big)\frac{B_h}{n}
     + \big(2H + \sqrt2\big) B_h \frac{\alpha^*}{n} \\
  &\le 2\sqrt2\sqrt{\Var_{\phat}(\Zt^t_{h+1})\frac{\alpha^*}{n}}
     + \frac1H \phat\abs{\Zt^t_{h+1} - Z^*_{h+1}}
     + 5 B_h \frac{\alpha}{n} + 4 H B_h \frac{\alpha}{n},
\end{align*}
since $\sqrt2 + 3 \le 5$ and $2H + \sqrt2 \le 4H$.
\end{proof}

\begin{remark}[The constants of the bonus]\label{rem:pos_constants_bonus}
\uses{def:bonus, lem:concentration_pos}
The proof of Lemma~7 in \cite{essakine2026tight} takes the square root of
$\Var_{p_h}(Z^*_{h+1})\alpha^*/n$ instead of $2\Var_{p_h}(Z^*_{h+1})\alpha^*/n$ and thereby
loses a factor $\sqrt2$ on the last two terms; the constants $2\sqrt2$, $5$ and $4H$ of its
statement are nevertheless correct, by the splitting $\sqrt{\alpha\alpha^*} \le (\alpha + \alpha^*)/2$
and the choice $\eta = 2/H$ above (the paper's $\eta = 1/H$ would give $\sqrt2 \cdot 4H$).
Two points about the bonus:
\begin{enumerate}
  \item[(i)] The third term of the bonus is stated here with the KL rate $\alpha(n,\delta)$,
        as in the statement of Lemma~7 of the paper (its equation (bonus positive) writes
        $\alpha^*$ there instead). Lemma~\ref{lem:concentration_pos} with $\alpha^*$ in the
        third term is true as well (the proof above gives the bound with $(2H + \sqrt2)\alpha^*$ in
        that place), but the certificate bound of Lemma~\ref{lem:certificate_pos} needs
        $b \ge (5 + 4H) B_h \alpha/n$, and this fails for the version with $\alpha^*$ (the
        KL-Bernstein inequality that controls $(p_h - \phat)(Z^*_{h+1} - \Zr^t_{h+1})$ only has the
        KL radius $\alpha/n$ at its disposal, and $\alpha^*(n,\delta)$ can be as small as
        $\alpha(n,\delta)/S$). \textbf{Consequently Definition~\ref{def:bonus} must use
        $4 H e^{\beta(H-h)}\alpha(n,\delta)/n$ as its third term} (for both signs of $\beta$);
        the sample complexity is unaffected, since the analysis of
        Lemma~\ref{lem:cert_true_recursion_pos} bounds this term by the additive rate term
        $84 H c_h \rho$, which is already in $\alpha$.
  \item[(ii)] The bonus is \emph{defined} as the right-hand side of
        Lemma~\ref{lem:concentration_pos} minus its last term, so the lemma is the statement
        ``$\abs{(p_h - \phat)Z^*_{h+1}} \le b + \frac1H\phat\abs{\Zt^t_{h+1} - Z^*_{h+1}}$'', which is
        the only form used afterwards (Lemmas~\ref{lem:optimism_pos} and~\ref{lem:certificate_pos}).
        The statement carries no sign assumption on $\Zt^t_{h+1} - Z^*_{h+1}$ (deviation~13);
        the absolute value is removed in Lemma~\ref{lem:optimism_pos} by its induction hypothesis.
\end{enumerate}
\end{remark}

\begin{lemma}[Optimism and pessimism; Lemma~8 of the paper]\label{lem:optimism_pos}
\lean{Essakine2026Tight.optExpValue_mem_Icc_optZ_of_pos, Essakine2026Tight.optExpQ_mem_Icc_stepUAt_of_pos, Essakine2026Tight.optExpQ_le_stepUAt_fst_of_pos, Essakine2026Tight.stepUAt_snd_le_optExpQ_of_pos, Essakine2026Tight.optExpValue_mem_Icc_ZlowerAt_ZtildeAt_of_pos, Essakine2026Tight.optExpQ_mem_Icc_UlowerAt_UtildeAt_of_pos, Essakine2026Tight.optExpValue_mem_Icc_of_pos, Essakine2026Tight.optZ_mem_of_pos}\leanok
\uses{def:good_event, def:backups, def:opt_value, def:entropic_bpi}
On $\evGood^+$, for all $t \ge 0$, $h \in \{1, \dots, H + 1\}$, $s \in \Ss$ and $a \in \As$,
\[
  \Ul_h^t(s,a) \le U^*_h(s,a) \le \Ut_h^t(s,a) \quad (h \le H)
  \qquad\text{and}\qquad
  \Zl_h^t(s) \le Z^*_h(s) \le \Zt_h^t(s) .
\]
\end{lemma}

\begin{proof}
\uses{lem:concentration_pos, lem:opt_exp_value_range, lem:opt_bellman, lem:emp_trans_simplex, def:backups, lem:backups_range}
\leanok
Fix $\omega \in \evGood^+$ and $t$. We prove by backward induction on $h$ (from $H + 1$ down to
$1$) the statement $P(h)$: ``for all $s$, $\Zl_h^t(s) \le Z^*_h(s) \le \Zt_h^t(s)$'', proving the
$U$-inequalities at step $h \le H$ along the way from $P(h+1)$.

\emph{Base case} $h = H + 1$: $\Zl_{H+1}^t = Z^*_{H+1} = \Zt_{H+1}^t = 1$
(Definitions~\ref{def:backups}, \ref{def:opt_value}).

\emph{Induction step.} Let $h \le H$, assume $P(h+1)$ and fix $(s,a)$. By
Lemma~\ref{lem:opt_bellman}, $U^*_h(s,a) = e^{\beta r_h(s,a)} (p_h Z^*_{h+1})(s,a)$ and
$U^*_h(s,a) \in [1, c_h]$ (Lemma~\ref{lem:opt_exp_value_range}).

\emph{Never-visited pair} ($n = 0$): $\Ut_h^t(s,a) = c_h \ge U^*_h(s,a) \ge 1 = \Ul_h^t(s,a)$.

\emph{Visited pair} ($n \ge 1$). Upper bound: $\Ut_h^t(s,a) = \min\{c_h, X\}$ with
$X := e^{\beta r_h(s,a)}[\phat\Zt^t_{h+1} + b + \frac1H\phat(\Zt^t_{h+1} - \Zl^t_{h+1})]$, and
$U^*_h(s,a) \le c_h$, so it suffices to show $U^*_h(s,a) \le X$. Now
\[
  X - U^*_h(s,a) = e^{\beta r_h(s,a)}\Big[ \phat(\Zt^t_{h+1} - Z^*_{h+1}) + (\phat - p_h)Z^*_{h+1}(s,a)
  + b + \tfrac1H \phat(\Zt^t_{h+1} - \Zl^t_{h+1}) \Big].
\]
By $P(h+1)$, $\Zt^t_{h+1} - Z^*_{h+1} \ge 0$ pointwise, so
$\phat\abs{\Zt^t_{h+1} - Z^*_{h+1}} = \phat(\Zt^t_{h+1} - Z^*_{h+1})$ and
Lemma~\ref{lem:concentration_pos} gives
$(\phat - p_h)Z^*_{h+1}(s,a) \ge -b - \frac1H\phat(\Zt^t_{h+1} - Z^*_{h+1})$. Hence
\[
  X - U^*_h(s,a) \ge e^{\beta r_h(s,a)}\Big[\big(1 - \tfrac1H\big)\phat(\Zt^t_{h+1} - Z^*_{h+1})
  + \tfrac1H \phat(Z^*_{h+1} - \Zl^t_{h+1})\Big] \ge 0 ,
\]
because $H \ge 1$, both functions inside are nonnegative by $P(h+1)$, and $\phat$ is a
positive linear functional (Lemma~\ref{lem:emp_trans_simplex}). Lower bound:
$\Ul_h^t(s,a) = \max\{1, Y\}$ with
$Y := e^{\beta r_h(s,a)}[\phat\Zl^t_{h+1} - b - \frac1H\phat(\Zt^t_{h+1} - \Zl^t_{h+1})]$; since
$U^*_h(s,a) \ge 1$ it suffices to show $Y \le U^*_h(s,a)$:
\[
  U^*_h(s,a) - Y = e^{\beta r_h(s,a)}\Big[\phat(Z^*_{h+1} - \Zl^t_{h+1}) + (p_h - \phat)Z^*_{h+1}(s,a)
  + b + \tfrac1H\phat(\Zt^t_{h+1} - \Zl^t_{h+1})\Big]
  \ge e^{\beta r_h(s,a)}\big(1 + \tfrac1H\big)\phat(Z^*_{h+1} - \Zl^t_{h+1}) \ge 0,
\]
using again Lemma~\ref{lem:concentration_pos} in the form
$(p_h - \phat)Z^*_{h+1}(s,a) \ge -b - \frac1H\phat(\Zt^t_{h+1} - Z^*_{h+1})$ and
$-\frac1H\phat(\Zt^t_{h+1} - Z^*_{h+1}) + \frac1H\phat(\Zt^t_{h+1} - \Zl^t_{h+1}) = \frac1H\phat(Z^*_{h+1} - \Zl^t_{h+1})$.

\emph{From $U$ to $Z$.} By Lemma~\ref{lem:opt_bellman} ($\beta > 0$), $Z^*_h(s) = \max_a U^*_h(s,a)$;
by Definition~\ref{def:backups}, $\Zt_h^t(s) = \max_a\Ut_h^t(s,a)$ and $\Zl_h^t(s) = \max_a \Ul_h^t(s,a)$;
the maximum over $a$ is monotone, which gives $P(h)$.
\end{proof}

\textbf{Lean remark.} The induction is \texttt{Fin.reverseInduction} on the Lean step
\texttt{h : Fin (H + 1)} (the values at \texttt{h} are \texttt{optZ r β δ hist (H - h)}), with the statement $P$
about $Z$ only, in the form $Z^*_h(s) \in [\Zl^t_h(s), \Zt^t_h(s)]$
(\texttt{optExpValue\_mem\_Icc\_optZ\_of\_pos}); the $U$-inequalities are separate lemmas taking
$P(h+1)$ as a hypothesis (\texttt{optExpQ\_le\_stepUAt\_fst\_of\_pos} for the upper bound,
\texttt{stepUAt\_snd\_le\_optExpQ\_of\_pos} for the lower bound), and the unconditional $U$-statement
is \texttt{optExpQ\_mem\_Icc\_stepUAt\_of\_pos}. The maxima over $a$ are LML's \texttt{Function.max}
(\texttt{optZ\_eq}), handled with \texttt{Function.le\_max} and \texttt{argmax\_spec}. The positivity of
$\phat$ is \texttt{vecExp\_mono}.

\section{The ring value and the certificate}
\label{sec:pos_certificate}

The pessimistic value $\Zl^t$ is below $Z^*$ but need not be comparable to $Z^{\pi}$
($\pi = \pi^{t+1}$ is greedy for $\Ut^t$, not for $\Ul^t$). The analysis-only value $\Zr^t$ of
Definition~\ref{def:ring} follows the policy $\pi$ under the \emph{true} kernel, clipped by the
pessimistic empirical backup; it lies below both $\Zl^t$ and $Z^\pi$, and the certificate
$\pi\cert^t$ dominates $\Zt^t - \Zr^t$. We first record the trivial lower bound on $\Zr^t$,
which the paper uses implicitly.

\begin{lemma}[The ring value is at least one]\label{lem:pos_ring_range}
\lean{Essakine2026Tight.one_le_ringZ_of_pos, Essakine2026Tight.one_le_ringU_of_pos, Essakine2026Tight.one_le_ringUAux_of_pos, Essakine2026Tight.one_le_ringZAt_of_pos, Essakine2026Tight.one_le_ringUAt_of_pos}\leanok
\uses{def:ring, def:entropic_bpi}
For all $t \ge 0$, $h \in \{1, \dots, H + 1\}$, $s$ and $a$: $\Zr_h^t(s) \ge 1$ and, for $h \le H$,
$\Ur_h^t(s,a) \ge 1$.
\end{lemma}

\begin{proof}
\uses{def:ring, lem:exp_value_monotone}
\leanok
Backward induction on $h$; $\Zr_{H+1}^t = 1$. If $\Zr_{h+1}^t \ge 1$ pointwise then
$e^{\beta r_h(s,a)}(p_h\Zr^t_{h+1})(s,a) \ge (p_h \Zr^t_{h+1})(s,a) \ge (p_h 1)(s,a) = 1$ ($\beta r_h \ge 0$;
Lemma~\ref{lem:exp_value_monotone}: $p_h$ is positive and $p_h 1 = 1$). For a visited pair,
$\Ur^t_{h,\mathrm{pes}}(s,a) = \max\{1, \cdot\} \ge 1$, so $\Ur_h^t(s,a)$, the minimum of two numbers
$\ge 1$, is $\ge 1$; for a never-visited pair $\Ur_h^t(s,a) = \min\{e^{\beta r_h}(p_h\Zr^t_{h+1})(s,a), 1\} \ge 1$.
Finally $\Zr_h^t(s) = \Ur_h^t(s, \pi_h(s)) \ge 1$.
\end{proof}

\begin{lemma}[The ring value is a lower bound; Lemma~10 of the paper]\label{lem:ring_pos}
\lean{Essakine2026Tight.ringZ_le_optZ_snd_and_expValue_of_pos, Essakine2026Tight.ringU_le_stepUAt_snd_and_expQ_of_pos, Essakine2026Tight.ringU_le_stepUAt_snd_of_pos, Essakine2026Tight.ringU_le_expQ_of_pos, Essakine2026Tight.ringZAt_le_ZlowerAt_and_expValue_of_pos, Essakine2026Tight.ringUAt_le_UlowerAt_and_expQ_of_pos}\leanok
\uses{def:ring, def:backups, def:greedy, def:exp_value, def:entropic_bpi}
For all $t \ge 0$ (no event is needed), $h \in \{1, \dots, H + 1\}$, $s$ and $a$, with $\pi = \pi^{t+1}$,
\[
  \Ur_h^t(s,a) \le \min\{\Ul_h^t(s,a), U^\pi_h(s,a)\} \quad (h \le H)
  \qquad\text{and}\qquad
  \Zr_h^t(s) \le \min\{\Zl_h^t(s), Z^\pi_h(s)\} .
\]
\end{lemma}

\begin{proof}
\uses{def:ring, def:backups, def:greedy, lem:exp_value_monotone, lem:exp_bellman, def:exp_value, lem:emp_trans_simplex}
\leanok
Backward induction on $h$ with the statement $P(h)$: ``for all $s$,
$\Zr_h^t(s) \le \min\{\Zl_h^t(s), Z^\pi_h(s)\}$''. Base case: all three values are $1$ at $h = H + 1$.
Let $h \le H$, assume $P(h+1)$, fix $(s,a)$.

\emph{Comparison with $U^\pi$.} In both the visited and the never-visited case
$\Ur_h^t(s,a) \le e^{\beta r_h(s,a)}(p_h\Zr^t_{h+1})(s,a)$ (it is a minimum with this quantity), and
$(p_h\Zr^t_{h+1})(s,a) \le (p_h Z^\pi_{h+1})(s,a)$ by $P(h+1)$ and the monotonicity of $p_h$
(Lemma~\ref{lem:exp_value_monotone}); by Definition~\ref{def:exp_value},
$e^{\beta r_h(s,a)}(p_h Z^\pi_{h+1})(s,a) = U^\pi_h(s,a)$.

\emph{Comparison with $\Ul$.} If $n = 0$, $\Ur_h^t(s,a) \le 1 = \Ul_h^t(s,a)$. If $n \ge 1$,
$\Ur_h^t(s,a) \le \Ur^t_{h,\mathrm{pes}}(s,a) = \max\{1, Y'\}$ and $\Ul_h^t(s,a) = \max\{1, Y\}$ with
\[
  Y = e^{\beta r_h}\big[\phat\Zl^t_{h+1} - b - \tfrac1H\phat(\Zt^t_{h+1} - \Zl^t_{h+1})\big], \qquad
  Y' = e^{\beta r_h}\big[\phat\Zr^t_{h+1} - b - \tfrac1H\phat(\Zt^t_{h+1} - \Zr^t_{h+1})\big],
\]
so $Y - Y' = e^{\beta r_h}\big[\phat(\Zl^t_{h+1} - \Zr^t_{h+1}) + \tfrac1H\phat(\Zl^t_{h+1} - \Zr^t_{h+1})\big]
= e^{\beta r_h}\big(1 + \tfrac1H\big)\phat(\Zl^t_{h+1} - \Zr^t_{h+1}) \ge 0$ by $P(h+1)$ and the
positivity of $\phat$ (Lemma~\ref{lem:emp_trans_simplex}); $y \mapsto \max\{1, y\}$ is monotone,
hence $\Ur^t_{h,\mathrm{pes}}(s,a) \le \Ul_h^t(s,a)$. (The paper writes the factor as $1 - \frac1H$; the
correct factor is $1 + \frac1H$, and only its sign matters.)

\emph{From $U$ to $Z$.} With $a = \pi_h(s)$: $\Zr_h^t(s) = \Ur_h^t(s, a) \le \Ul_h^t(s, a) \le \max_{a'}\Ul_h^t(s,a') = \Zl_h^t(s)$
(Definitions~\ref{def:ring}, \ref{def:backups}) and
$\Zr_h^t(s) \le U^\pi_h(s, \pi_h(s)) = Z^\pi_h(s)$ (Lemma~\ref{lem:exp_bellman}). This is $P(h)$.
\end{proof}

\begin{lemma}[One-step certificate bound; Lemma~11 of the paper]\label{lem:certificate_pos}
\lean{Essakine2026Tight.stepUAt_fst_sub_ringU_le_of_pos, Essakine2026Tight.UtildeAt_sub_ringUAt_le_of_pos}\leanok
\uses{def:good_event, def:ring, def:backups, def:bonus, def:entropic_bpi}
On $\evGood^+$, for all $t \ge 0$, $h \in \{1, \dots, H\}$ and $(s,a)$ with $n = n_h^t(s,a) \ge 1$,
\[
  \Ut_h^t(s,a) - \Ur_h^t(s,a)
  \le e^{\beta r_h(s,a)}\Big[3 b + \Big(1 + \frac3H\Big)\phat\big(\Zt^t_{h+1} - \Zr^t_{h+1}\big)\Big] .
\]
\end{lemma}

\begin{proof}
\uses{lem:concentration_pos, lem:optimism_pos, lem:ring_pos, lem:pos_ring_range, lem:kl_bernstein, lem:kl_transport_var, lem:var_le_range_mean, lem:sqrt_add_le, lem:sqrt_mul_le, lem:emp_trans_simplex, lem:backups_range, lem:opt_exp_value_range, def:event_kl, def:bonus, lem:pos_core_third_term}
\leanok
Fix $\omega \in \evGood^+$, $t$, $h$, $(s,a)$ with $n \ge 1$, and write $\Delta := \phat(\Zt^t_{h+1} - \Zr^t_{h+1}) \ge 0$
(Lemmas~\ref{lem:ring_pos}, \ref{lem:optimism_pos}: $\Zr^t_{h+1} \le \Zl^t_{h+1} \le Z^*_{h+1} \le \Zt^t_{h+1}$
pointwise). Throughout, $\Ut_h^t(s,a) \le X := e^{\beta r_h}[\phat\Zt^t_{h+1} + b + \frac1H\phat(\Zt^t_{h+1} - \Zl^t_{h+1})]$
(drop the clip), and $\phat(\Zt^t_{h+1} - \Zl^t_{h+1}) \le \Delta$ (positivity of $\phat$,
Lemma~\ref{lem:emp_trans_simplex}). Recall $\Ur_h^t(s,a) = \min\{e^{\beta r_h}(p_h\Zr^t_{h+1})(s,a), \Ur^t_{h,\mathrm{pes}}(s,a)\}$.

\emph{Case 1: $\Ur_h^t(s,a) = e^{\beta r_h}(p_h\Zr^t_{h+1})(s,a)$.} Then
$\Ut_h^t(s,a) - \Ur_h^t(s,a) \le e^{\beta r_h}\big[b + \frac1H\Delta + \phat\Zt^t_{h+1} - p_h\Zr^t_{h+1}\big]$ and
\[
  \phat\Zt^t_{h+1} - p_h\Zr^t_{h+1}
  = \Delta + (\phat - p_h)Z^*_{h+1}(s,a) + (p_h - \phat)\big(Z^*_{h+1} - \Zr^t_{h+1}\big)(s,a) .
\]
\emph{Second term.} By Lemma~\ref{lem:concentration_pos} and $\abs{\Zt^t_{h+1} - Z^*_{h+1}} = \Zt^t_{h+1} - Z^*_{h+1} \le \Zt^t_{h+1} - \Zr^t_{h+1}$
pointwise, $(\phat - p_h)Z^*_{h+1}(s,a) \le b + \frac1H\Delta$.
\emph{Third term.} Let $f := Z^*_{h+1} - \Zr^t_{h+1}$; then $0 \le f \le B_h - 1 \le B_h$ pointwise
(Lemma~\ref{lem:opt_exp_value_range} and Lemma~\ref{lem:pos_ring_range}). Since $\omega \in \evKL$,
$\KL(\phat\,\|\,p_h(\cdot\mid s,a)) \le \alpha/n$, and Lemma~\ref{lem:kl_bernstein} (with $p = \phat$,
$q = p_h(\cdot \mid s,a)$, range $B_h$) gives
\[
  (p_h - \phat) f(s,a) \le \sqrt{2\Var_{p_h}(f)(s,a)\frac{\alpha}{n}} + \frac23 B_h\frac{\alpha}{n} .
\]
The variance is under the \emph{true} kernel (the second argument of the divergence); the
paper writes it under $\phat$, which is the wrong side. We transport it back:
Lemma~\ref{lem:kl_transport_var} (first inequality) gives $\Var_{p_h}(f) \le 2\Var_{\phat}(f) + 4B_h^2\alpha/n$
and Lemma~\ref{lem:var_le_range_mean} gives $\Var_{\phat}(f) \le B_h\,\phat f$. Hence, by
Lemma~\ref{lem:sqrt_add_le} and Lemma~\ref{lem:sqrt_mul_le} (with $x = \phat f$, $y = 4B_h\alpha/n$,
$\eta = 2/H$),
\[
  \sqrt{2\Var_{p_h}(f)\frac{\alpha}{n}}
  \le \sqrt{\phat f \cdot 4 B_h \frac{\alpha}{n}} + \sqrt{8 B_h^2\frac{\alpha^2}{n^2}}
  \le \frac1H \phat f + H B_h\frac{\alpha}{n} + 2\sqrt2 B_h \frac{\alpha}{n} ,
\]
and $\phat f \le \Delta$. Therefore
\[
  (p_h - \phat)f(s,a) \le \frac1H\Delta + \Big(H + 2\sqrt2 + \frac23\Big) B_h\frac{\alpha}{n}
  \le \frac1H\Delta + (4H + 5) B_h \frac{\alpha}{n} \le \frac1H \Delta + b ,
\]
since $H + 2\sqrt2 + \frac23 \le H + 3.5 \le 4H + 5$ and $b \ge 5B_h\alpha/n + 4HB_h\alpha/n$
(Definition~\ref{def:bonus} with the third term in $\alpha$, Remark~\ref{rem:pos_constants_bonus};
the square-root term of $b$ is nonnegative).
\emph{Sum.} $\phat\Zt^t_{h+1} - p_h\Zr^t_{h+1} \le 2b + (1 + \frac2H)\Delta$, hence
$\Ut_h^t(s,a) - \Ur_h^t(s,a) \le e^{\beta r_h}[3b + (1 + \frac3H)\Delta]$.

\emph{Case 2: $\Ur_h^t(s,a) = \Ur^t_{h,\mathrm{pes}}(s,a)$.} Then
$\Ur_h^t(s,a) \ge e^{\beta r_h}[\phat\Zr^t_{h+1} - b - \frac1H\Delta]$ (the maximum dominates its
second argument), so
\[
  \Ut_h^t(s,a) - \Ur_h^t(s,a) \le e^{\beta r_h}\Big[2b + \Delta + \tfrac1H\phat(\Zt^t_{h+1} - \Zl^t_{h+1}) + \tfrac1H\Delta\Big]
  \le e^{\beta r_h}\Big[2b + \big(1 + \tfrac2H\big)\Delta\Big] ,
\]
which is at most the claimed bound as $b \ge 0$ (Lemma~\ref{lem:backups_range}) and $\Delta \ge 0$.
No event is used in this case.
\end{proof}

\begin{remark}[The constants $3$ and $3/H$]\label{rem:pos_constants_certificate}
\uses{lem:certificate_pos}
The paper's proof of Lemma~11 bounds the third term by
$\frac1H\Delta + (2H + \frac23)B_h\alpha/n$ and then claims that this is at most $\frac1H\Delta + b$.
With the bonus of equation (bonus positive) of the paper this would require
$(2H + \frac23) B_h\alpha/n \le 5B_h\alpha/n + 4HB_h\alpha^*/n$, which is false in general
($\alpha^*/\alpha$ can be as small as $1/S$). With the third term of the bonus in $\alpha$
(Remark~\ref{rem:pos_constants_bonus}) the inequality $(H + 2\sqrt2 + \frac23) B_h \alpha/n \le b$
holds for every $H \ge 1$, and the constants $3$ and $3/H$ of the certificate
(Definition~\ref{def:certificate}) are correct. The coefficient $\frac1H$ in front of $\Delta$ is
what makes the certificate recursion $(1 + \frac3H)$-contracting up to the constant $e^3$
(Lemma~\ref{lem:cert_unrolled_pos}); it is bought at the price of the $HB_h\alpha/n$ term, which is
why the bonus needs an $H\alpha/n$ term.
\end{remark}

\begin{lemma}[The certificate dominates the optimistic--ring gap]\label{lem:cert_dominates_gap_pos}
\lean{Essakine2026Tight.optZ_fst_sub_ringZ_le_cert_of_pos, Essakine2026Tight.ZtildeAt_sub_ringZAt_le_certAt_of_pos}\leanok
\uses{def:good_event, def:certificate, def:backups, def:ring, def:greedy, def:entropic_bpi}
On $\evGood^+$, for all $t \ge 0$, $h \in \{1, \dots, H + 1\}$ and $s \in \Ss$, with $\pi = \pi^{t+1}$,
\[
  \Zt_h^t(s) - \Zr_h^t(s) \le (\pi_h\cert_h^t)(s) .
\]
\end{lemma}

\begin{proof}
\uses{lem:certificate_pos, lem:backups_range, lem:pos_ring_range, lem:emp_trans_simplex, def:certificate, def:greedy, def:ring, lem:certificate_range}
\leanok
Fix $\omega \in \evGood^+$ and $t$; backward induction on $h$. Base case $h = H + 1$: both sides
are $0$ ($\Zt_{H+1}^t = \Zr_{H+1}^t = 1$, $\cert_{H+1}^t = 0$). Induction step: let $h \le H$, $s \in \Ss$
and $a := \pi_h(s)$; by Definitions~\ref{def:greedy} and~\ref{def:ring},
$\Zt_h^t(s) = \Ut_h^t(s,a)$ and $\Zr_h^t(s) = \Ur_h^t(s,a)$. In all cases
$\Zt_h^t(s) - \Zr_h^t(s) \le c_h - 1 \le c_h$ by Lemma~\ref{lem:backups_range} ($\Zt_h^t \le c_h$) and
Lemma~\ref{lem:pos_ring_range} ($\Zr_h^t \ge 1$).

\emph{Never-visited pair} ($n_h^t(s,a) = 0$): $(\pi_h\cert_h^t)(s) = c_h$ (Definition~\ref{def:certificate}),
and we are done.

\emph{Visited pair} ($n \ge 1$): $(\pi_h\cert_h^t)(s) = \min\{c_h, e^{\beta r_h(s,a)}[3b + (1 + \frac3H)\phat(\pi_{h+1}\cert_{h+1}^t)]\}$;
the first argument was handled above, and for the second, Lemma~\ref{lem:certificate_pos} and the
induction hypothesis $\Zt^t_{h+1} - \Zr^t_{h+1} \le \pi_{h+1}\cert^t_{h+1}$ pointwise, integrated
against the probability vector $\phat$ (Lemma~\ref{lem:emp_trans_simplex}), give
$\Zt_h^t(s) - \Zr_h^t(s) \le e^{\beta r_h(s,a)}[3b + (1 + \frac3H)\phat(\pi_{h+1}\cert^t_{h+1})]$.
\end{proof}

\begin{lemma}[The certificate dominates the suboptimality gap; Lemma~9 of the paper]\label{lem:stopping_pos}
\lean{Essakine2026Tight.optExpValue_sub_expValue_le_cert_of_pos, Essakine2026Tight.optExpValue_sub_expValue_le_certAt_of_pos}\leanok
\uses{def:good_event, def:certificate, def:opt_value, def:exp_value, def:greedy, def:entropic_bpi}
On $\evGood^+$, for all $t \ge 0$, $h \in \{1, \dots, H + 1\}$ and $s \in \Ss$, with $\pi = \pi^{t+1}$,
\[
  Z^*_h(s) - Z^\pi_h(s) \le (\pi_h\cert_h^t)(s) .
\]
\end{lemma}

\begin{proof}
\uses{lem:optimism_pos, lem:ring_pos, lem:cert_dominates_gap_pos}
\leanok
$Z^*_h(s) \le \Zt_h^t(s)$ (Lemma~\ref{lem:optimism_pos}) and $\Zr_h^t(s) \le Z^\pi_h(s)$
(Lemma~\ref{lem:ring_pos}), so $Z^*_h(s) - Z^\pi_h(s) \le \Zt_h^t(s) - \Zr_h^t(s) \le (\pi_h\cert_h^t)(s)$
(Lemma~\ref{lem:cert_dominates_gap_pos}).
\end{proof}

\begin{lemma}[Soundness of the stopping rule]\label{lem:pac_at_stopping_pos}
\lean{Essakine2026Tight.optEntropicValue_sub_entropicValue_le_of_stopCond_of_pos, Essakine2026Tight.optEntropicValue_sub_entropicValue_greedyAt_le_of_pos}\leanok
\uses{def:good_event, def:stopping_rule, def:entropic_value, def:opt_value, def:greedy, def:entropic_bpi}
On $\evGood^+$, for every $t \ge 0$ such that the stopping condition of
Definition~\ref{def:stopping_rule} holds after $t$ episodes, i.e.\
$(\pi_1\cert_1^t)(s_1) \le \frac{e^{\beta\eps} - 1}{e^{\beta\eps}}\Zt_1^t(s_1)$ with $\pi = \pi^{t+1}$,
one has $V^*_1(s_1) - V^\pi_1(s_1) \le \eps$.
\end{lemma}

\begin{proof}
\uses{lem:cert_dominates_gap_pos, lem:ring_pos, lem:stopping_pos, lem:value_gap_log, lem:exp_value_range}
\leanok
Write $g := (\pi_1\cert_1^t)(s_1) \ge 0$. Multiplying the stopping condition by $e^{\beta\eps} > 0$
gives $e^{\beta\eps} g \le (e^{\beta\eps} - 1)\Zt_1^t(s_1)$, i.e.\ $g \le (e^{\beta\eps} - 1)(\Zt_1^t(s_1) - g)$.
By Lemma~\ref{lem:cert_dominates_gap_pos} and Lemma~\ref{lem:ring_pos},
$\Zt_1^t(s_1) - g \le \Zr_1^t(s_1) \le Z^\pi_1(s_1)$, and $e^{\beta\eps} - 1 > 0$, so
$g \le (e^{\beta\eps} - 1) Z^\pi_1(s_1)$. By Lemma~\ref{lem:stopping_pos},
$Z^*_1(s_1) - Z^\pi_1(s_1) \le g \le (e^{\beta\eps} - 1)Z^\pi_1(s_1)$, and $Z^\pi_1(s_1) > 0$
(Lemma~\ref{lem:exp_value_range}); Lemma~\ref{lem:value_gap_log} ($\beta > 0$) concludes.
\end{proof}

\textbf{Lean remark.} The stopping condition is \texttt{stopCond r β δ ε s₁ (histAt X Y t ω)};
this lemma is about the history of $t$ episodes for a fixed $t$ and does not mention $\tau$. The
translation from the exponential gap to the entropic gap is entirely in
Lemma~\ref{lem:value_gap_log}; here only real arithmetic is used.

\section{The certificate under the true kernel}
\label{sec:pos_true_kernel}

\begin{lemma}[Certificate recursion under the true kernel]\label{lem:cert_true_recursion_pos}
\lean{Essakine2026Tight.cert_le_of_pos, Essakine2026Tight.certAt_le_of_pos, Essakine2026Tight.expValue_mem_Icc_of_pos}\leanok
\uses{def:good_event, def:certificate, def:rate_min, def:exp_value, def:greedy, def:entropic_bpi}
On $\evGood^+$, for all $t \ge 0$, $h \in \{1, \dots, H\}$ and $s \in \Ss$, with $\pi = \pi^{t+1}$ and
$a := \pi_h(s)$,
\[
  (\pi_h\cert_h^t)(s) \le e^{\beta r_h(s,a)}\Big[
    36\sqrt{\Var_{p_h}(Z^\pi_{h+1})(s,a)\,\rho^{*t}_h(s,a)}
    + \Big(1 + \frac{13}{H}\Big)\big(p_h\,\pi_{h+1}\cert^t_{h+1}\big)(s,a)
    + 84\,H\,c_h\,\rho^t_h(s,a)\Big] ,
\]
where $\rho^t_h(s,a) = \min\{\alpha(n,\delta)/n, 1\}$ and $\rho^{*t}_h(s,a) = \min\{\alpha^*(n,\delta)/n, 1\}$
(both equal to $1$ if $n = n_h^t(s,a) = 0$; Definition~\ref{def:rate_min}).
\end{lemma}

\begin{proof}
\uses{lem:certificate_range, def:certificate, def:bonus, def:event_kl, lem:kl_bernstein, lem:var_le_range_mean, lem:kl_transport_var, lem:transport_var, lem:sqrt_add_le, lem:sqrt_mul_le, lem:rates_props, lem:opt_exp_value_eq, lem:optimism_pos, lem:ring_pos, lem:cert_dominates_gap_pos, lem:backups_range, lem:exp_value_range, lem:exp_value_monotone, lem:emp_trans_simplex, lem:pos_core_recursion}
\leanok
Fix $\omega \in \evGood^+$, $t$, $h$, $s$, $a = \pi_h(s)$; write $g := \pi_{h+1}\cert^t_{h+1} : \Ss \to [0, B_h]$
(Lemma~\ref{lem:certificate_range}, the clip at step $h + 1$ being $c_{h+1} = B_h$),
$\Delta_p := (p_h g)(s,a) \ge 0$ and $\rho, \rho^*$ as in the statement. All three terms in the
bracket are nonnegative and $e^{\beta r_h} \ge 1$.

\emph{Trivial case: $n = 0$ or $\alpha/n \ge 1$.} Then $\rho = 1$, and $(\pi_h\cert_h^t)(s) \le c_h \le 84 H c_h \rho$
(Lemma~\ref{lem:certificate_range}), which is at most the right-hand side.

\emph{Main case: $n \ge 1$ and $\alpha/n < 1$.} Then $\rho = \alpha/n$ and, as $\alpha^* \le \alpha$
(Lemma~\ref{lem:rates_props}), $\rho^* = \alpha^*/n \le \rho$. Dropping the clip in
Definition~\ref{def:certificate}, $(\pi_h\cert_h^t)(s) \le e^{\beta r_h(s,a)}[3b + (1 + \frac3H)\phat g]$.

\emph{Step (i): the empirical expectation of $g$.} Since $\omega \in \evKL$,
Lemma~\ref{lem:kl_bernstein} (range $B_h$) gives
$\abs{(\phat - p_h) g(s,a)} \le \sqrt{2\Var_{p_h}(g)(s,a)\rho} + \frac23 B_h\rho$, and
$\Var_{p_h}(g)(s,a) \le B_h\Delta_p$ (Lemma~\ref{lem:var_le_range_mean}). By
Lemma~\ref{lem:sqrt_mul_le} with $x = \Delta_p$, $y = 2B_h\rho$, $\eta = 2/H$:
$\sqrt{2B_h\Delta_p\rho} \le \frac1H\Delta_p + \frac{H}{2}B_h\rho$. Hence, as $\frac H2 + \frac23 \le 2H$,
\[
  \phat g \le \Big(1 + \frac1H\Big)\Delta_p + 2 H B_h \rho .
\]

\emph{Step (ii): the bonus under the true kernel.} Recall
$b = 2\sqrt2\sqrt{\Var_{\phat}(\Zt^t_{h+1})\rho^*} + 5B_h\rho + 4HB_h\rho$ (Definition~\ref{def:bonus},
Remark~\ref{rem:pos_constants_bonus}). Lemma~\ref{lem:kl_transport_var} (second inequality, with
$f = \Zt^t_{h+1} - 1 \in [0, B_h]$ by Lemma~\ref{lem:backups_range}) gives
$\Var_{\phat}(\Zt^t_{h+1}) \le 2\Var_{p_h}(\Zt^t_{h+1})(s,a) + 4B_h^2\rho$.
Lemma~\ref{lem:transport_var} (first part, with $f = \Zt^t_{h+1} - 1$ and $g' = Z^\pi_{h+1} - 1$, both in
$[0, B_h]$ by Lemmas~\ref{lem:backups_range} and~\ref{lem:exp_value_range}) gives
$\Var_{p_h}(\Zt^t_{h+1}) \le 2\Var_{p_h}(Z^\pi_{h+1}) + 2B_h\,p_h\abs{\Zt^t_{h+1} - Z^\pi_{h+1}}$. Now
$Z^\pi_{h+1} \le Z^*_{h+1} \le \Zt^t_{h+1}$ pointwise (Lemma~\ref{lem:opt_exp_value_eq} for $\beta > 0$ and
Lemma~\ref{lem:optimism_pos}), and
$\Zt^t_{h+1} - Z^\pi_{h+1} \le \Zt^t_{h+1} - \Zr^t_{h+1} \le g$ pointwise
(Lemmas~\ref{lem:ring_pos} and~\ref{lem:cert_dominates_gap_pos}); by the monotonicity of $p_h$
(Lemma~\ref{lem:exp_value_monotone}), $p_h\abs{\Zt^t_{h+1} - Z^\pi_{h+1}}(s,a) \le \Delta_p$. Therefore
\[
  \Var_{\phat}(\Zt^t_{h+1})\rho^* \le 4\Var_{p_h}(Z^\pi_{h+1})(s,a)\rho^* + 4B_h\Delta_p\rho^* + 4B_h^2\rho\rho^* ,
\]
and by Lemma~\ref{lem:sqrt_add_le} (twice), Lemma~\ref{lem:sqrt_mul_le} with $x = \Delta_p$,
$y = 4B_h\rho^*$ and $\eta = 1/(\sqrt2 H)$ (so that $\sqrt{xy} \le \frac{x}{2\sqrt2 H} + \sqrt2 H\cdot 2B_h\rho^*$),
and $\sqrt{\rho\rho^*} \le \rho$,
\[
  2\sqrt2\sqrt{\Var_{\phat}(\Zt^t_{h+1})\rho^*}
  \le 4\sqrt2\sqrt{\Var_{p_h}(Z^\pi_{h+1})(s,a)\rho^*} + \frac1H\Delta_p + 8HB_h\rho^* + 4\sqrt2 B_h\rho .
\]
With $\rho^* \le \rho$, $4\sqrt2 \le 6$ and $8H + 4\sqrt2 + 5 + 4H \le (12 + 4\sqrt2 + 5)H \le 23H$,
\[
  b \le 6\sqrt{\Var_{p_h}(Z^\pi_{h+1})(s,a)\rho^*} + \frac1H\Delta_p + 23 H B_h\rho .
\]

\emph{Step (iii): combination.} Using (i), (ii), $(1 + \frac3H)(1 + \frac1H) = 1 + \frac4H + \frac3{H^2} \le 1 + \frac7H$
and $(1 + \frac3H)\cdot 2H = 2H + 6 \le 8H$ (all for $H \ge 1$),
\begin{align*}
  3b + \Big(1 + \frac3H\Big)\phat g
  &\le 18\sqrt{\Var_{p_h}(Z^\pi_{h+1})(s,a)\rho^*} + \frac3H\Delta_p + 69HB_h\rho
     + \Big(1 + \frac7H\Big)\Delta_p + 8HB_h\rho \\
  &\le 36\sqrt{\Var_{p_h}(Z^\pi_{h+1})(s,a)\rho^*} + \Big(1 + \frac{10}{H}\Big)\Delta_p + 77 H B_h\rho ,
\end{align*}
which is at most the bracket of the statement since $B_h \le c_h$, $10 \le 13$ and $77 \le 84$.
\end{proof}

\begin{remark}[The constants $36$, $13$, $84$ and the additive term]\label{rem:pos_constants_recursion}
\uses{lem:cert_true_recursion_pos}
The paper's computation (Appendix~B.1, ``Sample complexity'') is looser than the one above: it
keeps the coefficient $\frac{2\sqrt2}{H}$ in front of $\Delta_p$ in the bound on $b$, which gives the
coefficient $1 + \frac{4 + 6\sqrt2}{H} + \frac{3}{H^2} \le 1 + \frac{15.5}{H}$ (not $1 + \frac{13}{H}$) after
multiplication by $3$, and it drops a $+9$ when simplifying $81H + (1 + \frac3H)3H$ to $84H$.
Choosing $\eta = 1/(\sqrt2 H)$ instead of $2/H$ in the AM--GM step of (ii), and the sharper
form of Lemma~\ref{lem:sqrt_mul_le} in (i), gives $1 + \frac{10}{H}$ and $77H$, so the paper's
constants $36$, $13$, $84$ are correct with room to spare. The additive term is written with
$c_h = e^{\beta(H+1-h)}$ rather than the paper's $e^{\beta(H-h)}$ (deviation~11): it must dominate
the certificate of a never-visited pair, which is its clip $c_h$, and this is the only place
where the clip of Definition~\ref{def:certificate} enters the sample complexity. The rates are
the truncated rates $\rho = \min\{\alpha/n, 1\}$ (the paper's $\alpha(n)/n \wedge 1$), with the value
$1$ for never-visited pairs, so that the statement has no $\alpha(0)/0$.
\end{remark}

\textbf{Lean remark.} For the Lean step \texttt{h : Fin H} (the paper's $h + 1$) the statement is
about \texttt{cert r β δ hist (H - h) s} (\texttt{certAt X Y r β δ t ω h.castSucc s}) and the certificate
at the next step \texttt{cert r β δ hist (H - 1 - h)} (\texttt{certAt X Y r β δ t ω h.succ}); the clip
$c_h$ is \texttt{Real.exp (β * (H - h : ℕ))}, the variance is
\texttt{vecVar (M.transVec h s a) (expValue M β π h.succ)} and the rate terms are
\texttt{rateMin (alphaStar …) n} and \texttt{rateMin (alphaKL …) n} (\texttt{starRateMinAt},
\texttt{klRateMinAt} for a run). The proof follows the three steps above: step (i) is
Lemma~\ref{lem:pos_core_recursion}(i) with range $B_h$, step (ii) is Lemma~\ref{lem:pos_core_recursion}(ii)
with $G = \Zt^t_{h+1}$, $Z = Z^\pi_{h+1}$ and $c = 1$, and step (iii) is Lemma~\ref{lem:pos_core_recursion}(iii)
with $B = B' = B_h$ and $C = c_h$ (the proof gives the constant $77 B_h \le 84 c_h$).

\begin{lemma}[Unrolling with the factor $1 + c/H$; the terms of the unrolled bound]\label{lem:pos_core_unroll}
\lean{Learning.MDP.Episodic.le_exp_mul_sum_integral_of_backward, Learning.MDP.Episodic.sum_integral_exp_mul_sqrt_vecVar_mul_le, Learning.MDP.Episodic.sum_integral_exp_mul_mul_le, Learning.MDP.Episodic.integrable_comp_obs_action}\leanok
\uses{def:exp_value, def:step_law, def:occupancy}
Let $\pi$ be a policy, $s_1 \in \Ss$, $\beta \in \R$, $r : [H] \times \Ss \times \As \to \R$,
$\E^\pi = \E^\pi_{(s_1, 1)}$ and $w_h := e^{\beta\sum_{i \le h} r_i(S_i, A_i)}$ for $h \in [H]$.
\begin{enumerate}
  \item[(i)] Let $c \ge 0$, $u_h \ge 0$ and $g_h : \Ss \to \R$ ($h \le H + 1$) with $g_{H+1} = 0$ and
  $g_h(s) \le e^{\beta r_h(s,\pi_h(s))}[u_h(s, \pi_h(s)) + (1 + \frac cH)(p_h g_{h+1})(s, \pi_h(s))]$
  for all $h \le H$ and $s \in \Ss$. Then $g_1(s_1) \le e^c\sum_{h=1}^H\E^\pi[w_h\,u_h(S_h, A_h)]$.
  \item[(ii)] If $\mathcal M$ has the reward function $r$ and $\rho_h \ge 0$, then
  \[
    \sum_{h=1}^H\E^\pi\Big[w_h\sqrt{\Var_{p_h}(Z^\pi_{h+1})(S_h, A_h)\,\rho_h(S_h, A_h)}\Big]
    \le \sqrt{\Var_{\Prob^\pi}(e^{\beta R_1^\pi})}\,\sqrt{\sum_{h=1}^H\sum_{s,a}p^\pi_h(s,a)\rho_h(s,a)} .
  \]
  \item[(iii)] If $\rho_h \ge 0$ and $w_h\,c_h(S_h, A_h) \le C$ on every trajectory, then
  $\sum_h\E^\pi[w_h\,c_h(S_h, A_h)\,\rho_h(S_h, A_h)] \le C\sum_h\sum_{s,a}p^\pi_h(s,a)\rho_h(s,a)$.
\end{enumerate}
\end{lemma}

\begin{proof}
\uses{lem:unroll_recursion, lem:one_add_div_pow_le_exp, lem:cauchy_schwarz_traj, lem:entropic_variance_sum, lem:occupancy_expectation}
\leanok
(i) Lemma~\ref{lem:unroll_recursion} with $\kappa = 1 + c/H \ge 1$, and $\kappa^{h-1} \le e^c$ for
$1 \le h \le H$ (Lemma~\ref{lem:one_add_div_pow_le_exp}); the expectations are nonnegative.
(ii) Lemma~\ref{lem:cauchy_schwarz_traj} with the weights $w_h$, $v_h := \Var_{p_h}(Z^\pi_{h+1})(S_h, A_h) \ge 0$
and $u_h := \rho_h(S_h, A_h) \ge 0$; since $w_h^2 = e^{2\beta\sum_{i \le h} r_i(S_i,A_i)}$, the first factor is
$\Var_{\Prob^\pi}(e^{\beta R_1^\pi})$ by Lemma~\ref{lem:entropic_variance_sum}, and the second is the occupancy
sum by Lemma~\ref{lem:occupancy_expectation}. (iii) Integrate the pointwise bound
$w_h\,c_h\,\rho_h \le C\rho_h$ and apply Lemma~\ref{lem:occupancy_expectation}.
\end{proof}

\textbf{Lean remark.} These statements are MDP-general
(\texttt{LeanMachineLearning/ReinforcementLearning/MDP/UnrollBounds.lean}); they are recorded here
because they are the parts of Lemmas~\ref{lem:cert_unrolled_pos}, \ref{lem:ratio_bound_pos},
\ref{lem:cert_unrolled_neg} and~\ref{lem:ratio_bound_neg} that do not depend on the algorithm, and
would fit in Chapter~\ref{chap:pre_mdp}. The integrability side conditions are
\texttt{integrable\_exp\_sum\_mul} and \texttt{integrable\_comp\_obs\_action} (bounded functions of finitely
many rounds).

\begin{lemma}[Unrolled certificate]\label{lem:cert_unrolled_pos}
\lean{Essakine2026Tight.cert_le_exp_mul_sum_integral_of_pos, Essakine2026Tight.certAt_le_exp_mul_sum_integral_of_pos}\leanok
\uses{def:good_event, def:certificate, def:rate_min, def:exp_value, def:step_law, def:greedy, def:entropic_bpi}
On $\evGood^+$, for all $t \ge 0$, with $\pi = \pi^{t+1}$ and $\E^\pi = \E^\pi_{(s_1, 1)}$,
\[
  (\pi_1\cert_1^t)(s_1) \le e^{13}\,\E^\pi\Bigg[\sum_{h=1}^H e^{\beta\sum_{i \le h} r_i(S_i, A_i)}
  \Big(36\sqrt{\Var_{p_h}(Z^\pi_{h+1})(S_h, A_h)\,\rho^{*t}_h(S_h, A_h)}
  + 84\,H\,c_h\,\rho^t_h(S_h, A_h)\Big)\Bigg] .
\]
\end{lemma}

\begin{proof}
\uses{lem:cert_true_recursion_pos, lem:unroll_recursion, lem:one_add_div_pow_le_exp, def:certificate, lem:pos_core_unroll}
\leanok
Fix $\omega \in \evGood^+$ and $t$. Apply Lemma~\ref{lem:unroll_recursion} to the policy $\pi$, the
functions $g_h := \pi_h\cert_h^t$ ($h \le H + 1$; $g_{H+1} = 0$ by Definition~\ref{def:certificate}), the
nonnegative functions
$u_h(s,a) := 36\sqrt{\Var_{p_h}(Z^\pi_{h+1})(s,a)\rho^{*t}_h(s,a)} + 84Hc_h\rho^t_h(s,a)$ and
$\kappa := 1 + \frac{13}{H} \ge 1$: the hypothesis
$g_h(s) \le e^{\beta r_h(s, \pi_h(s))}[u_h(s, \pi_h(s)) + \kappa(p_hg_{h+1})(s, \pi_h(s))]$ for all
$h \le H$ and $s$ is Lemma~\ref{lem:cert_true_recursion_pos}, and the rewards are nonnegative. The
conclusion is $g_1(s_1) \le \sum_{h=1}^H\kappa^{h-1}\E^\pi[e^{\beta\sum_{i \le h} r_i(S_i,A_i)}u_h(S_h,A_h)]$, and
$\kappa^{h-1} \le \kappa^H \le e^{13}$ by Lemma~\ref{lem:one_add_div_pow_le_exp} (with $a = 13$,
$0 \le h - 1 \le H$); all the summands are nonnegative.
\end{proof}

\textbf{Lean remark.} The expectation is an integral against
$\mathrm{stepLaw}\ M\ \pi\ (s_1, 0)$ of a bounded measurable function of the first $H$ rounds;
the rate terms $\rho^t_h(S_h, A_h)$ are, for the fixed $\omega$, deterministic functions of $(h, s, a)$
(the history is frozen), so no measurability issue in $\omega$ arises inside the expectation.
The proof is Lemma~\ref{lem:pos_core_unroll}(i) with $c = 13$ and $g_h$ the certificate. For the
$0$-indexed Lean step $h$ the clip is \texttt{Real.exp (β * (H - h : ℕ))} ($= c_{h+1}$ in the
$1$-indexed notation above) and the weight is $e^{\beta\sum_{i \in \mathrm{Iic}\,h} r_i(S_i, A_i)}$.

\begin{lemma}[Normalized certificate bound]\label{lem:ratio_bound_pos}
\lean{Essakine2026Tight.cert_div_optZ_fst_le_of_pos, Essakine2026Tight.certAt_div_ZtildeAt_le_of_pos}\leanok
\uses{def:good_event, def:certificate, def:backups, def:rate_min, def:occupancy, def:gmax, def:upper_bound_const, def:greedy, def:entropic_bpi}
Let $V_G := (e^{\beta\Gmax} - 1)^2/e^{\beta\Gmax}$ with $\Gmax = \Gmax(\mathcal M)$ (Definition~\ref{def:gmax};
this is the $V_G$ of Definition~\ref{def:upper_bound_const}). On $\evGood^+$, for all $t \ge 0$, with
$\pi = \pi^{t+1}$,
\[
  \frac{(\pi_1\cert_1^t)(s_1)}{\Zt_1^t(s_1)} \le e^{13}\Big[36\sqrt{V_G\,x_t} + 84\,H\,e^{\beta(H+1)}\,y_t\Big],
  \quad
  x_t := \sum_{h=1}^H\sum_{s,a} p^\pi_h(s,a)\rho^{*t}_h(s,a), \quad
  y_t := \sum_{h=1}^H\sum_{s,a} p^\pi_h(s,a)\rho^t_h(s,a),
\]
where $p^\pi_h$ is the occupancy measure of $\pi$ from $s_1$ (Definition~\ref{def:occupancy}).
\end{lemma}

\begin{proof}
\uses{lem:cert_unrolled_pos, lem:optimism_pos, lem:opt_exp_value_eq, lem:exp_value_range, lem:cauchy_schwarz_traj, lem:entropic_variance_sum, lem:variance_ratio_le, lem:gmax_le_horizon, lem:occupancy_expectation, def:entropic_variance, lem:pos_core_unroll}
\leanok
Fix $\omega \in \evGood^+$ and $t$. \emph{Denominator.} $\Zt_1^t(s_1) \ge Z^*_1(s_1) \ge Z^\pi_1(s_1) \ge 1 > 0$
(Lemma~\ref{lem:optimism_pos}, Lemma~\ref{lem:opt_exp_value_eq} for $\beta > 0$,
Lemma~\ref{lem:exp_value_range}), so the left-hand side is at most $(\pi_1\cert_1^t)(s_1)/Z^\pi_1(s_1)$,
which we bound by dividing Lemma~\ref{lem:cert_unrolled_pos} by $Z^\pi_1(s_1)$.

\emph{Variance term.} Put $w_h := e^{\beta\sum_{i \le h} r_i(S_i,A_i)}$ (a nonnegative function of the
trajectory up to step $h$), $v_h := \Var_{p_h}(Z^\pi_{h+1})(S_h,A_h)/Z^\pi_1(s_1)^2$ and
$u_h := \rho^{*t}_h(S_h,A_h)$. Lemma~\ref{lem:cauchy_schwarz_traj} (Cauchy--Schwarz for $\sum_h\E^\pi$) gives
\[
  \E^\pi\Big[\sum_h w_h\sqrt{v_hu_h}\Big] \le \sqrt{\E^\pi\Big[\sum_h w_h^2v_h\Big]}\sqrt{\E^\pi\Big[\sum_h u_h\Big]} .
\]
By Lemma~\ref{lem:entropic_variance_sum},
$\E^\pi[\sum_h w_h^2 v_h] = \Var_{\Prob^\pi}(e^{\beta R_1^\pi})/Z^\pi_1(s_1)^2$, and by
Lemma~\ref{lem:variance_ratio_le} with $G = \Gmax$ (admissible since $R_1^\pi \in [0, \Gmax]$ a.s.\ by
Lemma~\ref{lem:gmax_le_horizon}), this ratio is at most $V_G$. By Lemma~\ref{lem:occupancy_expectation},
$\E^\pi[\sum_h u_h] = x_t$.

\emph{Rate term.} Since $0 \le r_i \le 1$ and $\beta > 0$, $w_h \le e^{\beta h}$, and $e^{\beta h}c_h = e^{\beta(H+1)}$;
with $Z^\pi_1(s_1) \ge 1$ and Lemma~\ref{lem:occupancy_expectation},
$\E^\pi[\sum_h w_h\,84Hc_h\rho^t_h(S_h,A_h)]/Z^\pi_1(s_1) \le 84He^{\beta(H+1)}\,y_t$.
\end{proof}

\textbf{Lean remark.} Lemma~\ref{lem:cauchy_schwarz_traj} allows the weight $w_h$ to depend on
the whole trajectory up to step $h$ (here on the past rewards), not only on $(S_h, A_h)$, which
is needed here. The identity of Lemma~\ref{lem:entropic_variance_sum} is where the entropic
variance $\sigma V^\pi_1(s_1)$ of Definition~\ref{def:entropic_variance} is identified with the
variance of $e^{\beta R_1^\pi}$ under $\Prob^\pi_{(s_1,1)}$. The variance and rate terms are
Lemma~\ref{lem:pos_core_unroll}(ii) and~(iii) (with $c_h = 84He^{\beta(H+1-h)}$ and
$C = 84He^{\beta(H+1)}$). The Lean statement writes $e^{|\beta|(H+1)}$ (equal to $e^{\beta(H+1)}$ here), so
that both signs produce the same hypothesis for Lemma~\ref{lem:pos_core_stopping_time}; the
ratio uses the greedy policy of the history, and the occupancy sums are
$\sum_{h, s, a}\texttt{occupancy M π s₁ h s a} \cdot \rho$.

\section{Bounding the stopping time}
\label{sec:pos_stopping_time}

\begin{lemma}[Progress at a non-stopping episode]\label{lem:progress_pos}
\lean{Essakine2026Tight.progress_le_cert_div_optZ_fst_of_pos, Essakine2026Tight.progress_le_certAt_div_ZtildeAt_of_pos}\leanok
\uses{def:good_event, def:stopping_rule, def:certificate, def:backups, def:upper_bound_const, def:entropic_bpi}
Let $\kappa(\beta, \eps) := (e^{|\beta|\eps} - 1)e^{-2|\beta|\eps}/2$ (Definition~\ref{def:upper_bound_const}).
On $\evGood^+$, for every $t \ge 0$ such that the stopping condition of Definition~\ref{def:stopping_rule}
\emph{fails} after $t$ episodes,
\[
  \kappa(\beta, \eps) \le \frac{(\pi_1\cert_1^t)(s_1)}{\Zt_1^t(s_1)} .
\]
\end{lemma}

\begin{proof}
\uses{def:stopping_rule, lem:backups_range}
\leanok
Failure of the condition means $(\pi_1\cert_1^t)(s_1) > \frac{e^{\beta\eps} - 1}{e^{\beta\eps}}\Zt_1^t(s_1)$, and
$\Zt_1^t(s_1) \ge 1 > 0$ (Lemma~\ref{lem:backups_range}), so the ratio exceeds
$(e^{\beta\eps} - 1)e^{-\beta\eps} \ge (e^{\beta\eps} - 1)e^{-2\beta\eps}/2 = \kappa(\beta,\eps)$ (as $e^{-\beta\eps} \le 1$ and $\frac12 \le 1$;
here $|\beta| = \beta$). No event is actually needed.
\end{proof}

\textbf{Lean remark.} No event is assumed; $\eps \ge 0$ and $\delta \in (0, 1]$ suffice ($\delta$ enters
through the range $\Zt_1^t(s_1) \ge 1$ of Lemma~\ref{lem:backups_range}).

\begin{lemma}[Core of the bound on the stopping time]\label{lem:pos_core_stopping_time}
\lean{Essakine2026Tight.natCast_le_upperBound, Essakine2026Tight.natCast_le_upperBound_of_mul_progress_le, ENat.toENNReal_le_ofReal_of_forall_natCast_le, Essakine2026Tight.one_le_upperBound, Essakine2026Tight.one_le_coeffLin, Essakine2026Tight.progress_pos, Essakine2026Tight.progress_le_one, Essakine2026Tight.progress_mul_coeffSqrt, Essakine2026Tight.progress_mul_coeffLin, Essakine2026Tight.varianceFactor_nonneg, Essakine2026Tight.exp_thirteen_le}\leanok
\uses{def:upper_bound_const, def:rates}
Let $\beta \ne 0$, $\eps > 0$, $\delta \in (0, 1]$, integers $S, A \ge 1$, $H \ge 0$, $G \in \R$, $V_G$, $\kappa = \kappa(\beta, \eps)$ as in
Definition~\ref{def:upper_bound_const}.
\begin{enumerate}
  \item[(i)] Let $T \in \N$ and $x_t \ge 0$, $y_t \in \R$ ($t < T$) be such that
  $\kappa \le e^{13}[36\sqrt{V_Gx_t} + 84He^{|\beta|(H+1)}y_t]$ for every $t < T$,
  $\sum_{t<T}x_t \le 16SAH\,\alpha^*(T-1,\delta)\log(T+1)$ and $\sum_{t<T}y_t \le 16SAH\,\alpha(T-1,\delta)\log(T+1)$.
  Then $T \le \mathrm{upperBound}(S, A, H, \beta, \eps, \delta, G)$.
  \item[(ii)] If $\tau \in \N \cup \{\infty\}$, $U \in \R$ and every integer $T \le \tau$ satisfies $T \le U$, then
  $\tau \le U$; in particular $\tau < \infty$.
\end{enumerate}
\end{lemma}

\begin{proof}
\uses{lem:self_bounding, lem:rates_props}
\leanok
(i) This is Steps~1--3 of the proof of Lemma~\ref{lem:stopping_time_bound_pos}, with the sequences
$x_t$, $y_t$ abstracted. If $T = 0$ the claim is $0 \le \mathrm{upperBound}$, and $\mathrm{upperBound} \ge 1$
because $A_0 \ge 0$ ($3SAH/\delta \ge 3$, or $\log 0 = 0$ if $H = 0$), $C \ge 0$ and $D \ge 0$. If $T \ge 1$ and
$H = 0$, the counting hypotheses give $x_0 \le 0$ and the hypothesis at $t = 0$ gives $\kappa \le 0$,
contradicting $\kappa > 0$. If $T, H \ge 1$: with $L = \log(8eT)$, Steps~1--3 give
$T\kappa \le 144e^{13}\sqrt{V_GSAH}\sqrt{T(A_0L + L^2)} + 1344e^{13}e^{|\beta|(H+1)}H^2SA(A_0L + SL^2)$
(the nonnegativity of $y_t$ is not needed), and $e^{13} \le 3^{13} = 1594323$ gives
$144e^{13} \le 10^9$ and $1344e^{13} \le 10^{10}$. Since $\kappa C = 10^9\sqrt{V_GSAH}$ and
$\kappa D = 10^{10}e^{|\beta|(H+1)}H^2SA$, dividing by $\kappa > 0$ gives the hypothesis of
Lemma~\ref{lem:self_bounding} ($A = A_0 \ge 0$, $B = 1$, $E = S \ge 1$, $\alpha = 8e$, $C \ge 0$, $D \ge 1$ because
$\kappa \le 1$ and $H^2SA \ge 1$), whose conclusion is $T \le \mathrm{upperBound}$.
(ii) If $\tau = \infty$, the integer $T = \lfloor U\rfloor_+ + 1 \le \tau$ satisfies $T > U$; otherwise take $T = \tau$.
\end{proof}

\begin{lemma}[Bound on the stopping time]\label{lem:stopping_time_bound_pos}
\lean{Essakine2026Tight.stoppingTime_le_upperBound_of_pos}\leanok
\uses{def:good_event, def:entropic_bpi, def:upper_bound_const, def:gmax}
At every $\omega \in \evGood^+$ at which the policies played are the greedy policies
($X_t(\omega) = \pi^{t+1}$ for all $t$, which holds almost surely by Lemma~\ref{lem:entropic_bpi_run}),
the stopping time of the run satisfies
$\tau \le \mathrm{upperBound}(S, A, H, \beta, \eps, \delta, \Gmax(\mathcal M))$ (Definition~\ref{def:upper_bound_const});
in particular $\tau < \infty$.
\end{lemma}

\begin{proof}
\uses{lem:progress_pos, lem:ratio_bound_pos, lem:sum_counts_bound, lem:entropic_bpi_run, lem:self_bounding, def:rates, def:upper_bound_const, lem:rates_props, def:rate_min, lem:pos_core_stopping_time}
\leanok
Fix $\omega \in \evGood^+$ (in particular $\omega \in \evCnt$) and write $\kappa := \kappa(\beta,\eps) > 0$,
$A_0 := \log(3SAH/\delta) > 0$, $C := 10^9\sqrt{V_G SAH}/\kappa$ and $D := 10^{10}e^{\beta(H+1)}H^2SA/\kappa$ (the
constants of Definition~\ref{def:upper_bound_const}, with $|\beta| = \beta$).

\emph{Step 1: an inequality for every $T \le \tau$.} Let $T \ge 1$ be an integer with $T \le \tau$
(i.e.\ the stopping condition fails after $t$ episodes for every $t < T$: by
Lemma~\ref{lem:entropic_bpi_run}, $\tau$ is the first $t$ at which it holds). Summing
Lemma~\ref{lem:progress_pos} and Lemma~\ref{lem:ratio_bound_pos} over $t = 0, \dots, T - 1$,
\[
  T\kappa \le \sum_{t < T}\frac{(\pi_1\cert_1^t)(s_1)}{\Zt_1^t(s_1)}
  \le 36e^{13}\sqrt{V_G}\sum_{t < T}\sqrt{x_t} + 84e^{13}He^{\beta(H+1)}\sum_{t < T}y_t .
\]
By the Cauchy--Schwarz inequality, $\sum_{t < T}\sqrt{x_t} \le \sqrt{T}\sqrt{\sum_{t < T}x_t}$.

\emph{Step 2: the counting bounds.} $x_t$ and $y_t$ are the inner sums of
Lemma~\ref{lem:sum_counts_bound} (with the policy $\pi^{t+1}$, the counts $n_h^t$ of the run and the
rates $\alpha^*$, resp.\ $\alpha$), so on $\evCnt$,
$\sum_{t < T}x_t \le 16SAH\,\alpha^*(T-1,\delta)\log(T+1)$ and $\sum_{t < T}y_t \le 16SAH\,\alpha(T-1,\delta)\log(T+1)$.
Put $L := \log(8eT)$; for $T \ge 1$, $\log(T+1) \le L$ and, by Definition~\ref{def:rates},
$\alpha^*(T-1,\delta) = A_0 + L$ and $\alpha(T-1,\delta) = A_0 + SL$. Hence
\[
  \sum_{t < T}x_t \le 16SAH\,(A_0L + L^2), \qquad \sum_{t < T}y_t \le 16SAH\,(A_0L + SL^2) .
\]

\emph{Step 3: the self-bounding inequality.} Combining,
\[
  T \le \frac{144\,e^{13}\sqrt{V_G SAH}}{\kappa}\sqrt{T(A_0L + L^2)}
     + \frac{1344\,e^{13}e^{\beta(H+1)}H^2SA}{\kappa}\,(A_0L + SL^2) .
\]
Numerically $e^{13} \le 4.43\cdot 10^5$, so $144\,e^{13} \le 6.4\cdot 10^7 \le 10^9$ and
$1344\,e^{13} \le 6.0\cdot 10^8 \le 10^{10}$; both brackets are nonnegative, so
\[
  T \le C\sqrt{T(A_0L + L^2)} + D\,(A_0L + SL^2), \qquad L = \log(8eT) .
\]
This is the hypothesis of Lemma~\ref{lem:self_bounding} with $\tau := T$, $A := A_0$, $B := 1$,
$E := S$, $\alpha := 8e \ge e$ (so that $\log(\alpha\tau) = L$) and the constants $C \ge 0$, $D$: indeed
$E \ge B$, $A_0 > 0$, and $D \ge 1$ (because $\kappa \le 1/8$: $(u - 1)/u^2 \le 1/4$ for $u > 1$).
Its conclusion is
\[
  T \le C^2(A_0 + 1)C_1^2 + \big(D + 2\sqrt{D}C\big)(A_0 + S)C_1^2 + 1 = \mathrm{upperBound}(S,A,H,\beta,\eps,\delta,\Gmax),
\]
with the $C_1$ of Lemma~\ref{lem:self_bounding} for $\alpha = 8e$, $A + E = A_0 + S$ and the constants
$C$, $D$, which is the $C_1$ of Definition~\ref{def:upper_bound_const} (Remark~\ref{rem:pos_c1_constant}).

\emph{Step 4: conclusion.} If $\tau = \infty$, every integer $T \ge 1$ satisfies $T \le \tau$, so every
$T \ge 1$ is at most the finite number $\mathrm{upperBound}$: contradiction. Hence $\tau < \infty$; if
$\tau \ge 1$, take $T = \tau$ in Step~3; if $\tau = 0$ the bound holds since $\mathrm{upperBound} \ge 1$.
\end{proof}

\textbf{Lean remark.} $\tau$ is \texttt{IdentAlg.stoppingTime}, with values in \texttt{ℕ∞}; the
statement ``for every $t < \tau$ the stopping condition fails'' is the characterization of the
hitting time (\texttt{hittingAfter}, Lemma~\ref{lem:entropic_bpi_run}), valid also when $\tau = \infty$.
The proof is organized as: (a) a lemma ``for all $T \ge 1$ with $(T : \texttt{ℕ∞}) \le \tau$,
$T \le \mathrm{upperBound}$'' proved on $\evGood^+$ by Steps~1--3; (b) the deduction $\tau \le \mathrm{upperBound}$
by a case split on $\tau = \top$. The step $T\kappa \le \sum_t(\cdot)$ is a sum of $T$ inequalities each
valid at an $\omega$ of the event, the sums $\sum_{t<T}x_t$ are \texttt{Finset.sum (Finset.range T)},
and the final numeric inequalities only need $e^{13} \le 3^{13}$ (from \texttt{Real.exp\_one\_lt\_three}:
$144 \cdot 3^{13} \approx 2.3\cdot 10^8$ and $1344 \cdot 3^{13} \approx 2.1\cdot 10^9$). When
$\Gmax(\mathcal M) = 0$ (all returns vanish a.s.), $V_G = 0$ and $C = 0$, which Lemma~\ref{lem:self_bounding}
allows. Steps~1--3 and~4 are Lemma~\ref{lem:pos_core_stopping_time}(i) and~(ii), with
$x_t = \sum_{h,s,a}p^{X_t}_h(s,a)\rho^{*t}_h(s,a)$ and $y_t$ likewise: the counting bound of
Lemma~\ref{lem:sum_counts_bound} is about the policies $X_t$ played, the normalized certificate bound about
the greedy policies $\pi^{t+1}$, hence the hypothesis $X_t(\omega) = \pi^{t+1}$ in the statement (Lean:
\texttt{hX : ∀ t, X t ω = greedyAt X Y r β δ t ω}). The Lean hypotheses are $\delta \in (0, 1]$ and
$\eps > 0$; the conclusion is \texttt{(τ ω : ℝ≥0∞) ≤ ENNReal.ofReal (upperBound …)}, and no hypothesis that
$(X, Y)$ is a run is needed (the characterization of $\tau$ as a hitting time is pointwise).

\begin{remark}[The constant $C_1$]\label{rem:pos_c1_constant}
\uses{lem:self_bounding, def:upper_bound_const}
The paper's Lemma~29 states the self-bounding inequality with
$C_1 = \frac85\log(11\alpha^2(A + E)(C + D))$, which is false
(Remark~\ref{rem:pana_self_bounding_paper}). Chapter~\ref{chap:pre_analysis} proves it with
$C_1 = \frac85\log(11\alpha^2(A + E)(C + \sqrt D)^2)$, and with $(C + D)^2$ in place of
$(C + \sqrt D)^2$ when $D \ge 1$ (which holds here); Definition~\ref{def:upper_bound_const}
uses this last form with $\alpha = 8e$, $A + E = A_0 + S$, so that the two agree. For
$C + D \ge 1$ this $C_1$ is larger than the paper's by $\frac85\log(C + D)$. Nothing else in
this chapter depends on the form of $C_1$; the asymptotic form $\tilde O(\kappa^{-2}V_GSAH)$ of
the bound is the same in both cases.
\end{remark}

\begin{lemma}[From an event of probability at least $1 - \delta$ to the probability statements]\label{lem:pos_core_pac}
\lean{Essakine2026Tight.one_sub_le_measureReal_stoppingTime_le, Essakine2026Tight.measureReal_bad_le, Essakine2026Tight.stoppingTime_ne_top_of_le}\leanok
\uses{def:entropic_bpi, def:stopping_rule}
Let $(X, Y, \mathrm{out})$ be a run of Entropic-BPI (with any parameters) in any environment, $E \subseteq \Omega$ with
$\Prob(E^c) \le \delta$, $U \in \R$, and call $\omega$ \emph{greedy} if $X_t(\omega) = \pi^{t+1}$ for all $t$.
\begin{enumerate}
  \item[(i)] If $\tau(\omega) \le U$ for every greedy $\omega \in E$, then $\Prob(\tau \le U) \ge 1 - \delta$.
  \item[(ii)] If $\tau(\omega) < \infty$ for every greedy $\omega \in E$, and for every $\omega \in E$ and $t$ such that the
  stopping condition holds after $t$ episodes $V^*_1(s_1) - V^{\pi^{t+1}}_1(s_1) \le \eps$, then
  $\Prob(V^*_1(s_1) - V^{\mathrm{out}}_1(s_1) > \eps) \le \delta$.
\end{enumerate}
\end{lemma}

\begin{proof}
\uses{lem:entropic_bpi_run}
\leanok
Let $N$ be the null set of Lemma~\ref{lem:entropic_bpi_run}. (i) $\{\tau \le U\}^c \setminus N \subseteq E^c$, so
$\Prob(\{\tau \le U\}^c) \le \delta$ and $\Prob(\tau \le U) \ge 1 - \Prob(\{\tau \le U\}^c)$ (subadditivity, no
measurability needed). (ii) For $\omega \in E \setminus N$, $\tau(\omega) < \infty$, the stopping condition holds after
$\tau(\omega)$ episodes and $\mathrm{out}(\omega) = \pi^{\tau(\omega)+1}$, so $\omega$ is not in the bad event; hence the bad
event is contained in $E^c \cup N$.
\end{proof}

\begin{lemma}[PAC guarantee and sample complexity of Entropic-BPI for $\beta > 0$]\label{lem:pac_pos}
\lean{Essakine2026Tight.measureReal_bad_le_of_pos, Essakine2026Tight.one_sub_le_measureReal_stoppingTime_le_of_pos}\leanok
\uses{def:entropic_bpi, def:entropic_value, def:opt_value, def:upper_bound_const, def:gmax, def:good_event}
In the standing setting,
\[
  \Prob\big(V^*_1(s_1) - V^{\mathrm{out}}_1(s_1) > \eps\big) \le \delta
  \qquad\text{and}\qquad
  \Prob\big(\tau \le \mathrm{upperBound}(S,A,H,\beta,\eps,\delta,\Gmax(\mathcal M))\big) \ge 1 - \delta ,
\]
(on the event $\{\tau = \infty\}$, of probability at most $\delta$ by the second statement, the
output $\mathrm{out}$ is whatever policy the run assigns).
\end{lemma}

\begin{proof}
\uses{lem:stopping_time_bound_pos, lem:entropic_bpi_run, lem:pac_at_stopping_pos, lem:good_event, lem:pos_core_pac}
\leanok
By Lemma~\ref{lem:good_event}, $\Prob(\evGood^+) \ge 1 - \delta$, so it suffices to show that, up to a
$\Prob$-null set $N$, $\evGood^+ \setminus N$ is contained in both events. Let $N$ be the null set
outside of which the conclusions of Lemma~\ref{lem:entropic_bpi_run} hold (the actions are the
greedy policies of the histories, the stopping condition holds at $\tau$ when $\tau < \infty$, and
$\mathrm{out} = \pi^{\tau + 1}$ on $\{\tau < \infty\}$). For $\omega \in \evGood^+ \setminus N$:
Lemma~\ref{lem:stopping_time_bound_pos} gives $\tau(\omega) \le \mathrm{upperBound} < \infty$; then the
stopping condition holds after $t := \tau(\omega)$ episodes and $\mathrm{out}(\omega) = \pi^{t+1}$, so
Lemma~\ref{lem:pac_at_stopping_pos} gives $V^*_1(s_1) - V^{\mathrm{out}(\omega)}_1(s_1) \le \eps$.
\end{proof}

\textbf{Lean remark.} This is the run-level statement (\texttt{IdentAlg.IsRun}) that
Chapter~\ref{chap:sample_complexity} converts into the PAC property of the algorithm
(\texttt{IdentAlg.IsPAC} is about the law of the output, Definition~\ref{def:entropic_pac}) and into
Theorem~\ref{thm:upper_bound_complexity}. The measurability of the events $\{V^*_1 - V^{\mathrm{out}}_1 > \eps\}$
and $\{\tau \le \mathrm{upperBound}\}$ is immediate: the policy type is finite and $\tau$ is a stopping time.
The proof needs $\evGood^+$ only through its three defining properties, and the inclusion
$\evGood^+ \setminus N \subseteq \{\tau \le \mathrm{upperBound}\} \cap \{V^*_1 - V^{\mathrm{out}}_1 \le \eps\}$ is proved
pointwise. In Lean the two statements are two lemmas for a run in \texttt{statesEnv M s₁}, with $\delta \in (0, 1]$
and $\eps > 0$, each an instance of Lemma~\ref{lem:pos_core_pac} with $E = \evGood^+$ and
$\Prob((\evGood^+)^c) \le \delta$ (\texttt{measure\_compl\_goodEvent\_le}, Lemma~\ref{lem:good_event}); in fact no
measurability of the events is used (outer measure and complements).
